\documentclass[11pt]{article}
\usepackage{ap-lit}
\title{Unit 2: Poetry}
\author{Paul Wang}
\date{November 2026}
\docpart{Unit 2 \textperiodcentered\ Study Guide \textperiodcentered\ Poetry 36--45\% of MCQ}

\begin{document}
\maketitle

\begin{abstract}
Poetry is two or three of the five multiple-choice passages and the whole of the first essay. A poem compresses: every line break, sound, and image is a choice, and the questions ask what each choice does. This guide works from the outside in: the speaker and situation (section 1), structure and form (2), sound and meter (3), word choice and imagery (4), and comparison, the figures that carry most of a poem's meaning (5). Section 6 is a plan for Question 1, the Poetry Analysis essay. Short quotations come from poems published before 1929; other examples are written for this guide. Each section ends with check questions; the key is on the last page.
\end{abstract}

\scoringline{Check questions: \blank\ / \qtotal \qquad Target: 12+ \qquad Term I confuse: \blank[6em]}

\section{Speaker and situation}
\label{sec:1}

The \textbf{speaker} is the voice of the poem, never assumed to be the poet. Before anything else, establish the \textbf{situation}: who is speaking, to whom, where, when, and why now (the \textbf{occasion}). Many poems are addressed to someone: a lover, a reader, the dead, an idea. Addressing someone absent or something non-human (``Death, be not proud'') is \textbf{apostrophe}. A \textbf{dramatic monologue} is a poem spoken by a character to a silent listener, revealing more than the speaker intends.

\textbf{Tone} is the speaker's attitude, and it usually moves. Track it stanza by stanza and name the \textbf{shift}: from despair to hope, from argument to concession, from description to reflection. In a sonnet the shift is called the \textbf{volta}.

\begin{keyidea}[Four questions for the first reading]
Who is speaking? To whom? What is the situation? Where does the tone change, and to what?
\end{keyidea}

\begin{q}{3}{2.1 Speaker}{The speaker is a constructed voice; do not assume it is the poet.}
A poem is written in the first person by a man and spoken by an elderly woman remembering her wedding. The best term for the voice of the poem is
\chh{the poet}{the narrator's foil}{the speaker}{the audience}
\end{q}

\begin{q}{1}{2.1 Apostrophe}{Addressing an absent person or a non-human thing = apostrophe.}
``O Sleep, why do you visit fools and pass me by?'' The line is an example of
\chh{apostrophe}{enjambment}{dramatic irony}{onomatopoeia}
\end{q}

\section{Structure and form}
\label{sec:2}

\subsection{Lines and stanzas}
The \textbf{line} is the poem's basic unit. A line that ends with a pause in sense or punctuation is \textbf{end-stopped}; one that runs into the next is \textbf{enjambed}, creating momentum, suspense, or a double meaning at the break. A strong pause inside a line is a \textbf{caesura}. Stanzas group lines: couplet (2), tercet (3), quatrain (4).

\subsection{Forms}
\begin{terms}
Shakespearean sonnet & 14 lines of iambic pentameter: three quatrains and a couplet, rhymed \emph{abab cdcd efef gg}. The couplet often turns, sums up, or undercuts. \\
Petrarchan sonnet & An octave (\emph{abbaabba}) that poses a problem and a sestet that answers or turns; the volta usually comes at line 9. \\
Villanelle & 19 lines on two rhymes with two refrains repeated in a fixed pattern; suits obsession and return. \\
Rondeau & A short French form with a refrain taken from the opening words; suits a phrase that keeps coming back. \\
Ode & A formal, elevated address to a person, object, or idea. \\
Elegy & A poem of mourning, often moving from grief to some consolation. \\
Ballad & A narrative poem in short stanzas, often with a refrain. \\
Free verse & No regular meter or rhyme; structure comes from line breaks, repetition, and syntax. \\
\end{terms}

\begin{trap}
Naming the form earns nothing by itself. What counts is how the form shapes meaning: why the turn comes where it does, what the couplet does to the quatrains, what a refrain gains each time it returns.
\end{trap}

\begin{q}{4}{2.2 Enjambment}{The sentence runs past the line break, so the break holds a moment of suspense.}
``I would have told her everything, if she / had asked me once.'' The break after ``she'' is an example of
\chh{caesura}{a refrain}{end-stopping}{enjambment}
\end{q}

\begin{q}{2}{2.2 Sonnet}{In a Shakespearean sonnet the final couplet often turns, concludes, or recasts the quatrains.}
In a Shakespearean sonnet, the closing couplet most often
\ch{introduces a new rhyme scheme borrowed from the octave}{turns, concludes, or recasts the argument of the three quatrains}{repeats the first line word for word}{shifts from iambic to trochaic meter}
\end{q}

\begin{q}{3}{2.2 Refrain}{A refrain returns with new weight as the context around it changes.}
A poem's refrain is repeated at the end of each stanza, but its final appearance ends with an exclamation point instead of a period. The change most likely signals
\ch{a printer's error}{that the refrain now means its literal opposite}{a heightened intensity in the speaker's final use of the phrase}{a new speaker}
\end{q}

\section{Sound and meter}
\label{sec:3}

\subsection{Meter}
A \textbf{foot} is a unit of stressed (/) and unstressed (\(\times\)) syllables.
\begin{compare}{Foot}{Pattern}{Example}
Iamb & \(\times\ /\) & be\emph{HOLD}; ``When \emph{I} do \emph{count} the \emph{clock} that \emph{tells} the \emph{time}'' \\
Trochee & \(/\ \times\) & \emph{TY}-ger; ``\emph{Ty}ger \emph{Ty}ger, \emph{burn}ing \emph{bright}'' \\
Anapest & \(\times\ \times\ /\) & in-ter-\emph{VENE} \\
Dactyl & \(/\ \times\ \times\) & \emph{MUR}-mur-ing \\
Spondee & \(/\ /\) & \emph{BLACK NIGHT} \\
\end{compare}
\textbf{Iambic pentameter} (five iambs) is the default line of English verse, close to speech. What matters on the exam is \textbf{variation}: a trochee at the start of a line, a run of spondees, or an extra syllable pushes against the expected rhythm and draws attention to the words where it happens.

\subsection{Rhyme and sound}
\begin{terms}
Rhyme scheme & Letters for end rhymes (\emph{abab}). \textbf{Internal rhyme} falls within a line; \textbf{slant rhyme} is near rhyme (\emph{soul / all}), which can sound unsettled. \\
Alliteration & Repeated initial consonant sounds: ``\emph{b}are ruin'd choirs \ldots\ \emph{b}lack night.'' \\
Assonance / consonance & Repeated vowel sounds / repeated consonant sounds anywhere in the words. \\
Euphony / cacophony & Smooth, open sounds / harsh, clashing sounds. \\
Onomatopoeia & Words that imitate sounds (\emph{hiss}, \emph{murmur}). \\
\end{terms}

\begin{q}{2}{2.3 Meter}{Five unstressed--stressed pairs = iambic pentameter.}
``The \emph{sun} has \emph{set} be\emph{hind} the \emph{ris}ing \emph{hill}.'' The line is written in
\chh{trochaic tetrameter}{iambic pentameter}{anapestic trimeter}{free verse}
\end{q}

\begin{q}{4}{2.3 Metrical variation}{A strong stress where an unstressed syllable is expected draws attention to the word.}
In a poem of regular iambic pentameter, one line begins with two heavily stressed syllables: ``\emph{Stop, stop}, my heart.'' The variation chiefly serves to
\ch{show that the poet has abandoned meter}{introduce a new rhyme scheme}{create a smooth, relaxed rhythm}{emphasize the urgency of the command}
\end{q}

\begin{q}{1}{2.3 Sound}{Harsh, clashing consonants = cacophony, suiting violence.}
A description of a battle uses lines such as ``crack and clang of bracket, brick, and blade.'' The sound effect is best described as
\chh{cacophony}{euphony}{slant rhyme}{onomatopoeia only}
\end{q}

\section{Word choice and imagery}
\label{sec:4}

Poems use few words, so each carries more. Look for \textbf{connotation} (``home'' versus ``residence''), \textbf{ambiguity} and \textbf{puns} (a word with two live meanings), and \textbf{diction level} (formal, colloquial, archaic, technical) and especially any \emph{change} in it. \textbf{Syntax} matters too: poets invert normal word order to place a word at the end of a line or to delay it, and a poem's single long sentence can feel very different from a string of short ones.

\textbf{Imagery} appeals to the senses: sight, sound, touch, taste, smell, and motion. Images often form a \textbf{pattern} (light and dark; growth and decay; heat and cold) that carries the poem's argument even when the speaker never states it.

\begin{trap}
A poem's title, final word, and any word that appears twice are almost always worth a question.
\end{trap}

\begin{q}{3}{2.4 Ambiguity}{``Lie'' means both rest and deceive here; both meanings are active.}
In ``Here let them lie, for lies are all they told,'' the first ``lie'' most clearly works as
\ch{an archaic spelling}{a simile}{a pun, meaning both ``rest'' and ``deceive''}{onomatopoeia}
\end{q}

\begin{q}{1}{2.4 Imagery pattern}{A movement from bright images to cold, dark ones traces a decline the speaker does not state.}
A poem's images move from ``noon,'' ``gold,'' and ``open fields'' in the first stanza to ``frost,'' ``ash,'' and ``shuttered rooms'' in the last. The pattern most likely conveys
\ch{a decline from vitality to loss}{a celebration of winter sports}{a description of a single afternoon only}{the speaker's love of architecture}
\end{q}

\section{Comparison: figurative language}
\label{sec:5}

\begin{terms}
Simile / metaphor & Explicit (``like,'' ``as'') or implied comparison. Ask what transfers: Wordsworth's ``I wandered lonely as a cloud'' makes loneliness drifting and weightless rather than painful. \\
Extended metaphor & One comparison developed through several lines or the whole poem. \\
Conceit & An elaborate, surprising comparison argued with ingenuity, as in the Metaphysical poets (lovers as the legs of a drawing compass). \\
Personification & Human qualities given to non-human things. \\
Metonymy / synecdoche & Naming a thing by something associated with it (``the crown'' for the monarch) / a part for the whole (``hands'' for workers). \\
Allusion & A reference to the Bible, myth, history, or another text, importing its meaning. \\
Paradox & A statement that seems contradictory but reveals a truth (``Death, thou shalt die''). \textbf{Oxymoron} compresses a paradox into two words (``sweet sorrow''). \\
Symbol & An image that means more than itself without being stated as a comparison. \\
\end{terms}

\begin{keyidea}[What transfers]
For every comparison, write two words for the qualities that move from the second thing to the first. A love ``like a red, red rose'' is fresh and beautiful, but roses also fade and have thorns, and a good essay notices when a poem uses the second set.
\end{keyidea}

\begin{q}{4}{2.5 Conceit}{An elaborate, ingenious, extended comparison between very unlike things = conceit.}
A poet compares two parted lovers to the two legs of a drawing compass: one stays fixed while the other moves, and their link brings the traveler back. This is best described as
\chh{an oxymoron}{onomatopoeia}{a synecdoche}{a conceit}
\end{q}

\begin{q}{2}{2.5 Synecdoche}{A part (sails) stands for the whole (ships).}
``Fifty sails came into the harbor at dawn.'' The word ``sails'' is an example of
\chh{apostrophe}{synecdoche}{paradox}{allusion}
\end{q}

\begin{q}{3}{2.5 Paradox}{A seeming contradiction that reveals a truth = paradox.}
``I must be cruel only to be kind.'' The line is best described as
\chh{a simile}{an allusion}{a paradox}{an apostrophe}
\end{q}

\section{Question 1: the Poetry Analysis essay}
\label{sec:6}

The prompt gives one poem (occasionally two to compare, though recent exams have used one), a sentence of context, and a task: \emph{analyze how the poet uses poetic elements and techniques to convey} a complex attitude, relationship, or experience.

\begin{enumerate}[leftmargin=1.6em]
\item \textbf{Read the prompt}, then the poem twice: once for the situation, once with a pen.
\item \textbf{Mark the turn.} Where does the tone, addressee, tense, or image pattern change? Your line of reasoning will usually follow the poem's movement up to and through that turn.
\item \textbf{Pick evidence of different kinds}: one structural choice (form, line break, the couplet), one sound or rhythm choice where it truly matters, and two or three images or comparisons.
\item \textbf{Name the complexity} the prompt asks about, in the thesis.
\item \textbf{Write about effects, not labels}: ``The enjambment after `hold' leaves the speaker's hand suspended in the empty space of the line, as he is suspended between holding on and letting go.''
\end{enumerate}

\begin{keyidea}[Thesis frame for Q1]
In ``[Title],'' [poet] moves from [first attitude] to [second attitude], using [technique] and [technique] to suggest that [insight about the complexity].
\end{keyidea}

\begin{q}{1}{2.6 Q1 thesis}{Names the movement and the complexity, and makes a defensible claim.}
A prompt asks how a poet conveys ``the speaker's complex attitude toward a lost home.'' Which thesis best responds?
\ch{Moving from bright remembered rooms to a closing image of a locked gate, the speaker treats the lost home as both a refuge memory can enter and a place the present has barred for good.}{The poem uses imagery, rhyme, and metaphor to talk about a home.}{The speaker misses his home.}{The poem is a sonnet about a house the speaker once lived in.}
\end{q}

\begin{q}{4}{2.6 Effect over label}{Explains how the technique works and what it does to meaning.}
Which sentence best analyzes a technique, as the rubric rewards?
\ch{The poem has an ABAB rhyme scheme.}{Line 4 contains alliteration.}{The poet uses many metaphors.}{By breaking the line after ``ever,'' the poet leaves ``ever'' hanging alone, so the promise sounds endless before the next line limits it.}
\end{q}

\makekey
\end{document}
